% §5.5 The base, and §5.6 the wedge of an orbit --- blueprint chapter 5, part 05-orbits.tex (Q-0288).
% Included from 05-landscape.tex; the conventions are in that file's header. Order: the blueprint's.
% The stable rung comes first in 5.6 because prop:orbit-step-cone reads prop:stable-step-cone;
% prop:stable-order-generated (blueprint part 05-wedges) follows prop:stable-exponent-cone, whose cone
% it reads (Q-0290, review flag S12). The two Matern cross-shape nodes and the grid remark are in
% appendix B, three routine proofs in appendix E (Q-0290). The six [A] nodes all rest on A3(iii),
% named once at the start of 5.6 (the review's length judgement). The identification of the Matern
% orbit with the anisotropic Matern of spatial statistics is a
% naming read at the page (records/affine/READING-whittle1954stationary.md,
% READING-lindgren2011explicit.md), not a fact a proof rests on.

\subsection{The base: the general linear group along rays}\label{sec:affine-orbits}

The smallest order of the scheme is the ray order. It requires only that the record at
$\lambda^2\Sigma$, $\lambda > 1$, be a further smoothing of the record at $\Sigma$: a zoom out without
a change of shape. This subsection asks which kernels can then be carried to every elliptical aperture
covariantly under the whole general linear group, and the answer is that this costs nothing beyond
what the isotropic theory requires. The kernel $\mu_I$ at the round aperture fixes the family, which is
its \emph{affine orbit}: the kernel at $\Sigma$ is $(\Sigma^{1/2})_*\mu_I$, the round kernel stretched
by $\Sigma^{1/2}$, as the Gaussian of covariance $2b\Sigma$ is the round Gaussian of covariance $2bI$
stretched. The orbit is a family along the rays exactly when $\mu_I$ is rotation invariant and
self-decomposable \ref{eq:self-decomposable} (Proposition~\ref{prop:ray-families}). The
rotation-invariant self-decomposable laws other than $\delta_0$ are the kernels from scale zero of the
isotropic families (Theorem~\ref{thm:isotropic-classification} and
Lemma~\ref{lem:radial-marginal}(1)), so every isotropic scale space, the stable and the \Matern{} ones
among them, has an affine orbit along the rays. What a non-Gaussian orbit cannot do is cascade along a
step that widens the aperture in some directions and not in others, a step of rank below $d$
(Proposition~\ref{prop:affine-orbits}(3)). The \Matern{} orbit of
Proposition~\ref{prop:affine-orbits}(4) is the example to keep in mind.

% shared with blueprint prop:affine-orbits
\begin{proposition}[affine orbits of isotropic laws]\label{prop:affine-orbits}
Let $d \ge 1$ and let $g : [0,\infty) \to [0,\infty)$ be continuous with $g(0) = 0$ and such that
$\omega \mapsto e^{-g(|\omega|^2)}$ is the characteristic function of a probability measure
$\mu_I$ on $\RR^d$. For $\Sigma \in \PD(d)$ put $\mu_\Sigma := (\Sigma^{1/2})_*\mu_I$, the
\emph{affine orbit} of $\mu_I$.
\begin{enumerate}
\item \emph{(Covariance.)} $\widehat\mu_\Sigma(\omega) = e^{-g(\langle\omega,\Sigma\omega\rangle)}$,
and $A_*\mu_\Sigma = \mu_{A\Sigma A^{\tp}}$ for every $A \in \GL(d)$ and $\Sigma \in \PD(d)$.
\item \emph{(Cascade along rays.)} If for some $\lambda > 1$ there is a probability measure
$\nu_\lambda$ on $\RR^d$ with $e^{-g(\lambda^2|\omega|^2)} = e^{-g(|\omega|^2)}\,\widehat\nu_\lambda(\omega)$,
then $\mu_{\lambda^2\Sigma} = \mu_\Sigma * (\Sigma^{1/2})_*\nu_\lambda$ for every
$\Sigma \in \PD(d)$.
\item \emph{(No cascade across shapes.)} If $d \ge 2$ and $g$ is not linear, then for no
$\Sigma \in \PD(d)$ and no $P \in \PSD(d)$ with $0 < \rank P < d$ is
$\mu_{\Sigma+P} = \mu_\Sigma * \rho$ with $\rho$ a probability measure. In particular the orbit is
not the family of kernels from scale zero of any affine cascade family.
\item \emph{(The \Matern{} orbits.)} For $\gamma > 0$, $g_\gamma(u) := \gamma\log(1+u)$ satisfies
the hypothesis in every dimension: $(1+|\omega|^2)^{-\gamma}$ is the characteristic function of
the law of $\sqrt{2V}\,Z$, with $V$ Gamma distributed of shape $\gamma$ and scale $1$ and $Z$
standard normal on $\RR^d$, independent. Its orbit has
$\widehat\mu_\Sigma(\omega) = (1 + \langle\omega,\Sigma\omega\rangle)^{-\gamma}$; clause (3)
applies to it, $g_\gamma$ not being linear; and the hypothesis of clause (2) holds for every
$\gamma > 0$ and every $\lambda > 1$, with $\nu_\lambda$ the law of $\sqrt{2Y}\,Z$, where $Z$ is
standard normal on $\RR^d$ and $Y \ge 0$, independent of $Z$, is compound Poisson with \Levy{}
measure $\gamma\,(e^{-x/\lambda^2} - e^{-x})\,x^{-1}dx$ on $(0,\infty)$, of total mass
$2\gamma\log\lambda$.
\end{enumerate}
\end{proposition}

Clause (3) is Theorem~\ref{thm:rank-deficient-step} read on the orbit. The stable orbit
$e^{-\langle\omega,\Sigma\omega\rangle^{\alpha/2}}$, $0 < \alpha < 2$, is another instance of the
proposition (Lemma~\ref{lem:elliptical-stable}), so by clause (3) it is not a family along the Loewner
order. For $\gamma > d/2$ the family of clause (4) is the \Matern{} covariance of spatial statistics
(\sscite{matern1960spatial}, (2.4.5)--(2.4.7), pp.~17--18; \sscite{whittle1954stationary}, \S9,
pp.~447--448; \sscite{lindgren2011explicit}, \S2.1, eqs.~(1)--(2), p.~426), with geometric anisotropy
(\sscite{lindgren2011explicit}, \S3.4, p.~437), in which the matrix is a parameter of the covariance. Clause (3) excludes
the rank-deficient steps between two values of it, and Proposition~\ref{prop:matern-steps} shows that
no step across shapes has a positive increment.
The proof is in Appendix~\ref{app:proofs}.


The converse holds: a family covariant under the general linear group that cascades along the ray of
$I$ is such an orbit, and the law must be self-decomposable.

% shared with blueprint prop:ray-families+663841edc318
\begin{proposition}[the ray families are the orbits of the self-decomposable isotropic laws]\label{prop:ray-families}
Let $d \ge 1$ and let $(\mu_\Sigma)_{\Sigma\in\PD(d)}$ be probability measures on $\RR^d$ with
\begin{enumerate}
\item[(R1)] \emph{(affine covariance)} $A_*\mu_\Sigma = \mu_{A\Sigma A^{\tp}}$ for every
$A \in \GL(d)$ and $\Sigma \in \PD(d)$;
\item[(R2)] \emph{(a cascade along the ray of $I$)} for every $\lambda > 1$ there is a probability
measure $\rho_\lambda$ on $\RR^d$ with $\mu_{\lambda^2I} = \mu_I * \rho_\lambda$.
\end{enumerate}
Then $\mu_I$ is rotation invariant and self-decomposable \ref{eq:self-decomposable},
$\mu_\Sigma = (\Sigma^{1/2})_*\mu_I$ for every $\Sigma$, and the family cascades along every ray:
$\mu_{\lambda^2\Sigma} = \mu_\Sigma * (\Sigma^{1/2})_*\rho_\lambda$ for every $\Sigma \in \PD(d)$
and $\lambda > 1$. Conversely, for every rotation-invariant self-decomposable law $\mu_I$ on
$\RR^d$ the family $\mu_\Sigma := (\Sigma^{1/2})_*\mu_I$ satisfies (R1) and (R2).
\end{proposition}

With Theorem~\ref{thm:affine-rigidity} this places the obstruction for $d \ge 2$: along the rays every
rotation-invariant self-decomposable law keeps its orbit, and across every step of the Loewner order
only the Gaussian does. Which steps across shapes
an orbit survives depends on the law, and that is the next subsection.

% proof of prop:ray-families
\begin{proof}
(R1) is (K3) of Definition~\ref{def:affine-cascade-family} for the kernels from scale zero, and (R2)
is (K1) along the ray of $I$ alone.

\emph{Covariance determines the orbit.} (R1) with $A = Q \in \Orth(d)$ and $\Sigma = I$ gives
$Q_*\mu_I = \mu_{QQ^{\tp}} = \mu_I$, and with $A = \Sigma^{1/2}$ it gives
$(\Sigma^{1/2})_*\mu_I = \mu_\Sigma$, as for the orbits of Proposition~\ref{prop:affine-orbits}.

\emph{Self-decomposability.} By (R1) with $A = \lambda I$, $\mu_{\lambda^2I} = (\lambda I)_*\mu_I$,
so (R2) reads $(\lambda I)_*\mu_I = \mu_I * \rho_\lambda$. Applying $(\lambda^{-1}I)_*$ gives
$\mu_I = (cI)_*\mu_I * (cI)_*\rho_\lambda$ with $c := \lambda^{-1}$, which is
\ref{eq:self-decomposable} with the residual $(cI)_*\rho_\lambda$, and every $c \in (0,1)$ arises.

\emph{Every ray.} $\mu_{\lambda^2\Sigma} = (\Sigma^{1/2})_*\mu_{\lambda^2I}
= (\Sigma^{1/2})_*(\mu_I * \rho_\lambda) = \mu_\Sigma * (\Sigma^{1/2})_*\rho_\lambda$.

\emph{Converse.} For (R1), let $B := A\Sigma^{1/2}$ and $\Sigma' := A\Sigma A^{\tp} = BB^{\tp}$. Then
$Q := \Sigma'^{-1/2}B$ is orthogonal, so
$A_*\mu_\Sigma = B_*\mu_I = (\Sigma'^{1/2})_*Q_*\mu_I = \mu_{\Sigma'}$. For (R2), self-decomposability
at $c := \lambda^{-1}$ gives $\mu_I = (cI)_*\mu_I * \rho_c$, and applying $(\lambda I)_*$ gives
$\mu_{\lambda^2I} = \mu_I * (\lambda I)_*\rho_c$.
\end{proof}

\subsection{The step set of an orbit}\label{sec:affine-stable-cone}

Shape adaptation changes the shape of the aperture as the scale grows (\S\ref{sec:landscape-uses}, the third use).
For the Gaussian every Loewner step is a further smoothing, so the shape may change as fast as
Lemma~\ref{lem:loewner-scale-shape}(3) allows: in the plane, a step that doubles the volume scale
reaches an anisotropy gained of $4$ and no more. This subsection asks the same of the other affine
orbits. Along which steps $\Sigma \to \Sigma'$ across shapes is the kernel at $\Sigma'$ the kernel at
$\Sigma$ convolved with a probability measure? Along those steps the coarser record is a further
smoothing of the finer one, and along every other step a change of shape is a re-measurement of the
image.

To state the question for every orbit, relate two apertures $\Sigma, \Sigma' \in \PD(d)$
when $\mu_{\Sigma'} = \mu_\Sigma * \rho$ for a positive finite measure $\rho$. The orbit is covariant
under $\GL(d)$ (Proposition~\ref{prop:affine-orbits}(1)), so this relation is $\GL(d)$-invariant, and by
Proposition~\ref{prop:invariant-orders}(1) it is $\le_{\mathcal K}$ for exactly one step set
$\mathcal K$: the set of $M$ with $\mu_M = \mu_I * \rho$. Call $\mathcal K$ the \emph{step set} of the
orbit. It contains the rays exactly when $\mu_I$ is self-decomposable
(Proposition~\ref{prop:ray-families}). For a non-Gaussian orbit it contains no step
$\Sigma \to \Sigma + P$ with $0 < \rank P < d$ (Proposition~\ref{prop:affine-orbits}(3)), so the
question is about steps of full rank.

The answer depends on the kernel. The orbit of the isotropic $\alpha$-stable law, $0 < \alpha < 2$,
cascades on full-rank steps of every orientation, under a speed limit of the Gaussian's form and
narrower: with $t_1, t_2$ the volume scales of the two apertures, the anisotropy gained is at most
$(t_2/t_1)^{d\alpha/(d+\alpha)}$, against $(t_2/t_1)^d$ for the Gaussian
(Proposition~\ref{prop:stable-step-cone}(3) and Lemma~\ref{lem:loewner-scale-shape}(3)). For the
Poisson kernel, $\alpha = 1$ (\S\ref{sec:classical-prior-art}), a step in the plane that doubles the
volume scale reaches an anisotropy gained of $2^{2/3} \approx 1.59$ and no more, and
Figure~\ref{fig:shape-plane} draws both limits. At the other end, the \Matern{} orbits cascade along
the rays and along no other step (Proposition~\ref{prop:matern-steps}). No orbit whose transform falls
off like a power of the frequency, in the sense of Proposition~\ref{prop:log-orbit-rays}, has a
positive increment on a step off the rays, and the orbits of decaying rational symbols are among them
(Proposition~\ref{prop:rational-symbols}). For the orbit of a general profile, one function of the
profile, its dilation ratio, gives a set of steps on which the orbit cascades
(Proposition~\ref{prop:orbit-step-cone}); whether it cascades on more steps is open. The subsection
ends with what the theorem on the general linear row (\S\ref{sec:wedge-axis-proof}) takes from the
orbits: the exponents of the isotropic stable law, the bound they put on a non-Gaussian law, and the
stable order extended to $\PSD(d)$.

The answers are read through the \Levy--Khintchine representation (C2, Theorem~\ref{thm:cited-A3}),
and from here on the subsection uses its vocabulary. The kernels of the stable orbits and of the
orbits of Proposition~\ref{prop:orbit-step-cone} are \emph{infinitely divisible} laws with no
Gaussian part. Each is the law of a sum of independent displacements, and its \emph{\Levy{} measure}
gives the expected number of displacements of each size and direction. Here that measure has a
density $n_\Sigma(x)$, the rate of the displacements per unit volume, not the density of the kernel.
A step $\Sigma \to \Sigma'$ is a further smoothing when the kernel at $\Sigma'$ has at least the rate
of displacements of the kernel at $\Sigma$ at every size and in every direction,
$n_\Sigma \le n_{\Sigma'}$ (Proposition~\ref{prop:orbit-step-cone}(2)), and for the stable orbits only
then (Proposition~\ref{prop:stable-step-cone}(1)).

\subsubsection*{The stable rung}

The chain of cones $C_\alpha$ of \S\ref{sec:invariant-wedges} is a ladder from the rays ($\alpha = 0$)
to the Loewner order ($\alpha = \infty$). Call the \emph{rung} of an orbit the largest index $\alpha$
whose stable order its step set contains. The elliptical stable laws are built from their \Levy{}
densities.

% shared with blueprint lem:elliptical-stable
\begin{lemma}[the elliptical stable laws]\label{lem:elliptical-stable}
Let $d \ge 1$ and $0 < \alpha < 2$, and put
\[
  c_{d,\alpha} := \int_{\RR^d}(1 - \cos y_1)\,|y|^{-d-\alpha}\,dy .
\]
Then $0 < c_{d,\alpha} < \infty$. For $H \in \PD(d)$ and $x \in \RR^d \setminus \{0\}$ put
\begin{equation}\label{eq:stable-levy-density}
  n_H(x) := c_{d,\alpha}^{-1}(\det H)^{-1/2}\langle x,H^{-1}x\rangle^{-(d+\alpha)/2} ,
  \qquad \psi_H(\omega) := \langle\omega,H\omega\rangle^{\alpha/2} .
\end{equation}
\begin{enumerate}
\item For every $H \in \PD(d)$ and $\omega \in \RR^d$,
$\psi_H(\omega) = \int_{\RR^d\setminus\{0\}}(1 - \cos\langle\omega,x\rangle)\,n_H(x)\,dx$, and
$\int_{\RR^d\setminus\{0\}}(|x|^2 \wedge 1)\,n_H(x)\,dx < \infty$.
\item For every $H \in \PD(d)$ there is a unique probability measure $S_\alpha(H)$ on $\RR^d$ with
$\widehat{S_\alpha(H)} = e^{-\psi_H}$. It is infinitely divisible with generating triplet
$(0, n_H(x)\,dx, 0)$, and $A_*S_\alpha(H) = S_\alpha(AHA^{\tp})$ for every $A \in \GL(d)$.
\item $n_{AHA^{\tp}}(x) = |\det A|^{-1}\,n_H(A^{-1}x)$ for $A \in \GL(d)$, $H \in \PD(d)$ and
$x \ne 0$.
\end{enumerate}
In particular $g(u) := u^{\alpha/2}$ satisfies the hypothesis of Proposition~\ref{prop:affine-orbits}, with
$\mu_I = S_\alpha(I)$, and the affine orbit of $S_\alpha(I)$ is $\mu_\Sigma = S_\alpha(\Sigma)$, the
\emph{stable orbit}.
\end{lemma}

The law $S_\alpha(H)$ is the rotation-invariant $\alpha$-stable law (\sscite{sato1999levy}, Thm.~14.14,
p.~86) transported by $H^{1/2}$, with its constant normalized to $1$. In the terms of Samorodnitsky
and Taqqu it is a sub-Gaussian stable law. The existence in (2) is their Proposition~2.5.2 with
$R = 2H$, and their Proposition~2.5.7 makes every such law a linear image of $S_\alpha(I)$
(\sscite{samorodnitsky1994stable}, pp.~78 and~81). At $\alpha = 1$, $S_1(I)$ is the kernel of the Poisson scale
space at unit scale, with transform $e^{-|\omega|}$ (\S\ref{sec:classical-prior-art}). The proof is in
Appendix~\ref{app:proofs}.

The stable orbit cascades on a cone of full-rank steps, the steps whose \Levy{} densities
\ref{eq:stable-levy-density} are ordered.

% shared with blueprint prop:stable-step-cone+2382c75d66ad
\begin{proposition}[the stable orbit cascades on a cone of full-rank steps]\label{prop:stable-step-cone}
Let $d \ge 1$, $0 < \alpha < 2$, and let $\mu_\Sigma = S_\alpha(\Sigma)$, $\Sigma \in \PD(d)$, be the
stable orbit (Lemma~\ref{lem:elliptical-stable}). For $\Sigma, \Sigma' \in \PD(d)$ put
$M := \Sigma^{-1/2}\Sigma'\Sigma^{-1/2}$ and let $\lambda_{\min}(M)$ be its least eigenvalue.
\begin{enumerate}
\item \emph{(The step cone.)} The following are equivalent:
(a) there is a positive finite measure $\rho$ with $\mu_{\Sigma'} = \mu_\Sigma * \rho$;
(b) $n_\Sigma(x) \le n_{\Sigma'}(x)$ for every $x \ne 0$;
(c) $\lambda_{\min}(M)^{d+\alpha} \ge \det M$.
When they hold, $\Sigma \le \Sigma'$, and $\rho$ is unique: it is the infinitely divisible probability
measure $\rho_{\Sigma,\Sigma'}$ with generating triplet $(0, (n_{\Sigma'} - n_\Sigma)(x)\,dx, 0)$.
\item \emph{(The order.)} Write $\Sigma \preceq_\alpha \Sigma'$ when (c) holds. This is a partial order
on $\PD(d)$, contained in the Loewner order, and $\Sigma \preceq_\alpha \Sigma'$ if and only if
$A\Sigma A^{\tp} \preceq_\alpha A\Sigma'A^{\tp}$, for every $A \in \GL(d)$. The increments satisfy
$\rho_{\Sigma,\Sigma} = \delta_0$, $\rho_{\Sigma,\Sigma''} = \rho_{\Sigma',\Sigma''} * \rho_{\Sigma,\Sigma'}$
whenever $\Sigma \preceq_\alpha \Sigma' \preceq_\alpha \Sigma''$, and
$A_*\rho_{\Sigma,\Sigma'} = \rho_{A\Sigma A^{\tp},A\Sigma'A^{\tp}}$.
\item \emph{(Which steps.)} Every ray, $\Sigma' = \lambda^2\Sigma$ with $\lambda \ge 1$, satisfies
$\Sigma \preceq_\alpha \Sigma'$, and no step $\Sigma' = \Sigma + P$ with $P \in \PSD(d)$ and
$0 < \rank P < d$ does. In the coordinates of Lemma~\ref{lem:loewner-scale-shape} (volume scales $t_1, t_2$,
the eigenvalues $n_1 \le \dots \le n_d$ of the relative shape $N$),
\[
  \Sigma_1 \preceq_\alpha \Sigma_2 \Longleftrightarrow n_1 \ge (t_1/t_2)^{2\alpha/(d+\alpha)} ,
\]
so the anisotropy gained, $\sqrt{n_d/n_1}$, is at most $(t_2/t_1)^{d\alpha/(d+\alpha)}$; and for
every $t_2 \ge t_1$ and every $N \in \PD(d)$ with $\det N = 1$ and least eigenvalue at least
$(t_1/t_2)^{2\alpha/(d+\alpha)}$, the matrix $\Sigma_2 := t_2^2S_1^{1/2}NS_1^{1/2}$ satisfies
$\Sigma_1 \preceq_\alpha \Sigma_2$. In particular, for $d \ge 2$ and $t_2 > t_1$ the relative shape
of a step in the cone may be any $N \ne I$ in that range, with any eigenvectors.
\end{enumerate}
\end{proposition}

So the stable orbits, which are not Gaussian, cascade across shapes, on full-rank steps of every
orientation. What
limits a step is its relative anisotropy, the gap between $\lambda_{\min}(M)$ and the geometric mean
$(\det M)^{1/d}$, and not its orientation. The cone lies between the rays and the Loewner order: the
Loewner order is $n_1 \ge (t_1/t_2)^2$ (Lemma~\ref{lem:loewner-scale-shape}(1)). Clause (3)
excludes the steps of rank below $d$ by condition (c), as Theorem~\ref{thm:rank-deficient-step} does for
$g(u) = u^{\alpha/2}$. By Proposition~\ref{prop:stable-order-generated}(1) the order is that of the
cone field $C_\alpha$.

The proof of (a)$\Rightarrow$(b) rests on a lemma that reads the \Levy{} densities off the tail of an
increment. Positivity is tested against functions supported at a large scale, so only a weak form of
the tail is needed, and the homogeneity of the exponent gives it.

% shared with blueprint lem:stable-tail
\begin{lemma}[the tail of an increment is the difference of the \Levy{} densities]\label{lem:stable-tail}
Let $d \ge 1$, $0 < \alpha < 2$, $H, H' \in \PD(d)$, and $C$ a real $d \times d$ matrix.
Suppose $\rho$ is a positive finite measure on $\RR^d$ with
\begin{equation}\label{eq:stable-tail-transform}
  \widehat\rho(\omega) = \exp\bigl[-\langle\omega,C\omega\rangle - \psi_{H'}(\omega) + \psi_H(\omega)\bigr]
  \qquad (\omega \in \RR^d),
\end{equation}
with $\psi_H$ and $n_H$ as in \ref{eq:stable-levy-density}. Then $\rho$ is a probability measure,
and for every smooth $\phi$ with compact support in $\RR^d \setminus \{0\}$
\begin{equation}\label{eq:stable-tail}
  \lim_{R\to\infty} R^\alpha\int\phi(x/R)\,\rho(dx) = \int\phi(x)\,\bigl(n_{H'}(x) - n_H(x)\bigr)\,dx .
\end{equation}
Consequently $n_H(x) \le n_{H'}(x)$ for every $x \ne 0$.
\end{lemma}

Since $n_{H'} - n_H$ is homogeneous of degree $-d-\alpha$, the right-hand side of \ref{eq:stable-tail}
equals $R^\alpha\int\phi(x/R)(n_{H'} - n_H)(x)\,dx$ for every $R$: at large scales the increment's mass
is that of the signed difference of \Levy{} densities. The limit compares masses at large scales, not
densities. The Gaussian factor $\langle\omega,C\omega\rangle$ is there for
Proposition~\ref{prop:similarity-stable-order-kernels}(3).

The proof writes $\int\phi(x/R)\,\rho(dx)$ by Fourier inversion as an integral of $\widehat\phi$ against
$e^{-E(\cdot/R)} - 1$, with $\widehat\rho = e^{-E}$; the term $1$ costs nothing because $\phi$ vanishes
at the origin. As the non-Gaussian part of the exponent is homogeneous of degree $\alpha$ and its Gaussian part of
degree $2 > \alpha$, the factor $R^\alpha$ turns this into $\int\phi\,(n_{H'} - n_H)\,dx$ in the limit, and a nonnegative bump supported where
$n_{H'} < n_H$ would make it negative. The complete proof is in Appendix~\ref{app:wedge-axis}.

% proof of prop:stable-step-cone
\begin{proof}
\emph{(b)$\Leftrightarrow$(c).} Substitute $x = \Sigma^{1/2}y$. Then
$\langle x,\Sigma^{-1}x\rangle = |y|^2$, $\langle x,\Sigma'^{-1}x\rangle = \langle y,M^{-1}y\rangle$ and
$\det\Sigma'/\det\Sigma = \det M$, so (b) says $\langle y,M^{-1}y\rangle/|y|^2 \le
(\det M)^{-1/(d+\alpha)}$ for every $y \ne 0$. The supremum of the left-hand side is $1/\lambda_{\min}(M)$,
so (b) is (c).

\emph{(c) gives $\Sigma \le \Sigma'$.} $\det M \ge \lambda_{\min}(M)^d$, so (c) gives $\lambda_{\min}(M) \ge 1$,
hence $M \ge I$ and $\Sigma' - \Sigma = \Sigma^{1/2}(M - I)\Sigma^{1/2} \ge 0$.

\emph{(a)$\Rightarrow$(b).} $e^{-\psi_\Sigma}$ has no zero, so
$\widehat\rho = e^{-\psi_{\Sigma'} + \psi_\Sigma}$, which is \ref{eq:stable-tail-transform} with $C = 0$,
and Lemma~\ref{lem:stable-tail} gives (b). The same identity shows that $\rho$ is unique.

\emph{(b)$\Rightarrow$(a).} $f := n_{\Sigma'} - n_\Sigma$ is nonnegative, even and at most $n_{\Sigma'}$,
so by Lemma~\ref{lem:elliptical-stable}(1) and C2(iii) there is an infinitely divisible $\rho$ with
triplet $(0, f(x)\,dx, 0)$ and $\widehat\rho = e^{-\psi_{\Sigma'} + \psi_\Sigma}$. So
$\mu_\Sigma * \rho$ and $\mu_{\Sigma'}$ have the same transform.

\emph{Clause (2).} Transitivity is that of (b), antisymmetry that of the Loewner order, and by
Lemma~\ref{lem:elliptical-stable}(3) condition (b) for $(A\Sigma A^{\tp}, A\Sigma'A^{\tp})$ is (b) for
$(\Sigma, \Sigma')$ at $A^{-1}x$. The transforms $e^{-\psi_{\Sigma'} + \psi_\Sigma}$ telescope, which
gives the cascade, and $\psi_\Sigma(A^{\tp}\omega) = \psi_{A\Sigma A^{\tp}}(\omega)$ gives the
covariance.

\emph{Clause (3).} A ray has $M = \lambda^2I$, and (c) reads $\lambda^{2(d+\alpha)} \ge \lambda^{2d}$. For
$0 < \rank P < d$, $M = I + \Sigma^{-1/2}P\Sigma^{-1/2}$ has the eigenvalue $1$ and one above $1$, so
$\lambda_{\min}(M) = 1 < \det M$. In the coordinates of Lemma~\ref{lem:loewner-scale-shape},
$M = (t_2/t_1)^2N$, so (c) is $n_1^{d+\alpha} \ge (t_1/t_2)^{2\alpha}$, and
$n_d/n_1 \le n_1^{-d} \le (t_2/t_1)^{2d\alpha/(d+\alpha)}$. The converse is
Lemma~\ref{lem:loewner-scale-shape}(3). For the last sentence take
$N = Q\diag(a, a^{-1/(d-1)}, \dots, a^{-1/(d-1)})Q^{\tp}$ with $Q$ arbitrary and
$(t_1/t_2)^{2\alpha/(d+\alpha)} \le a < 1$.
\end{proof}

\begin{remark}[a second route to (a)$\Rightarrow$(b)]\label{rem:stable-step-sato}
The implication (a)$\Rightarrow$(b) of Proposition~\ref{prop:stable-step-cone}(1) also follows from the
uniqueness of the generating triplet. The increment $\rho$ has the transform $e^{-E}$ with
$E := \psi_{\Sigma'} - \psi_\Sigma$, which is homogeneous of degree $\alpha$. So the image of $\rho$
under $x \mapsto n^{-1/\alpha}x$ is an $n$-th convolution root of $\rho$ for every $n$, and $\rho$ is
infinitely divisible. The generating triplets of $\mu_\Sigma$ and $\rho$ then add to one of
$\mu_{\Sigma'}$. By the uniqueness of the triplet (\sscite{sato1999levy}, Thm.~8.1(ii), pp.~37--38) the
\Levy{} measure of $\rho$ is $(n_{\Sigma'} - n_\Sigma)(x)\,dx$. So $n_{\Sigma'} - n_\Sigma \ge 0$
almost everywhere, and by continuity at every $x \ne 0$.
\end{remark}

The theorem on the general linear row needs a bound on how anisotropic the exponents of a
rotation-invariant non-Gaussian law can be. A law is \emph{operator-self-decomposable} with exponent
$B$ when it is its own image under $e^{-tB}$ convolved with a residual, for every $t > 0$
(\S\ref{sec:cited-facts}). Applying $(e^{tB})_*$, this says that the law stretched by $e^{tB}$ is the
law convolved with a probability measure, a further smoothing. For $B = I$ the stretch is a zoom, and
the condition is self-decomposability \ref{eq:self-decomposable}. On an affine orbit the kernel at $I$
stretched by $e^{tB}$ is the kernel at $e^{tB}e^{tB^{\tp}}$ (Proposition~\ref{prop:affine-orbits}(1)),
so an exponent of $\mu_I$ is a direction in which the orbit leaves the round aperture along further
smoothings. The stable law gives the bound, and for a single linear flow $e^{tB}$ the same computation
gives the set of its exponents.

% shared with blueprint prop:stable-exponent-cone+9fbc223a4e90
\begin{proposition}[the exponent cone of the isotropic stable law]\label{prop:stable-exponent-cone}
Let $d \ge 1$ and $B$ a real $d \times d$ matrix, and write $\mathrm{sym}\,B := \tfrac12(B + B^{\tp})$.
Call a probability measure $\mu$ on $\RR^d$ \emph{$B$-self-decomposable} if for every $t > 0$ there is a
probability measure $\rho_t$ with $\mu = (e^{-tB})_*\mu * \rho_t$. For $0 < \alpha < 2$, $S_\alpha(I)$
is $B$-self-decomposable if and only if
\begin{equation}\label{eq:stable-exponent-cone}
  (d+\alpha)\,\lambda_{\min}(\mathrm{sym}\,B) \ge \tr B .
\end{equation}
The set $\mathfrak k_\alpha$ of such $B$ is a closed convex cone, invariant under $B \mapsto QBQ^{\tp}$ for
$Q \in \Orth(d)$; it contains every skew-symmetric matrix and $I$, and every $B \in \mathfrak k_\alpha$ has
$\mathrm{sym}\,B \ge 0$, with $\lambda_{\min}(\mathrm{sym}\,B) > 0$ unless $B$ is skew-symmetric.
\end{proposition}

Only $\mathrm{sym}\,B$ enters, so the skew part of $B$, the turning of the flow, is free; what is
limited is the anisotropy of $\mathrm{sym}\,B$. $B = I$ is ordinary self-decomposability. The cone grows
with $\alpha$. The closure of the union of the $\mathfrak k_\alpha$ is the cone that
\ref{eq:stable-exponent-cone} defines at $\alpha = 2$, and for $d \ge 2$ it is smaller than the
Gaussian's cone $\{\mathrm{sym}\,B \ge 0\}$.
The proof is in Appendix~\ref{app:proofs}.

The exponent cones of the stable laws bound the contracting exponents of every rotation-invariant
non-Gaussian law: such an exponent lies in $\mathfrak k_\alpha$ for some $\alpha < 2$. The proof of
Theorem~\ref{thm:wedge-axis} uses this bound.

% shared with blueprint prop:anisotropy-budget
\begin{proposition}[the anisotropy budget of a rotation-invariant law]\label{prop:anisotropy-budget}
Let $d \ge 1$, let $B$ be a real $d \times d$ matrix, and write
$\lambda := \lambda_{\min}(\mathrm{sym}\,B)$ (Proposition~\ref{prop:stable-exponent-cone}).
\begin{enumerate}
\item If $e^{-tB} \to 0$ as $t \to \infty$ and $\mu$ is a rotation-invariant $B$-self-decomposable
law on $\RR^d$ that is not Gaussian (a point mass counting as Gaussian), then $\lambda > 0$ and
\[
  \tr B < (d+2)\,\lambda .
\]
\item Conversely, if $\tr B < (d+2)\lambda$, then $\lambda > 0$, $e^{-tB} \to 0$, and $S_\alpha(I)$ is
rotation invariant, not Gaussian, and $B$-self-decomposable for every $\alpha \in (0,2)$ with
$\alpha \ge \tr B/\lambda - d$.
\end{enumerate}
\end{proposition}

The ratio $\tr B/\lambda$ is at least $d$, with equality for a multiple of $I$ plus a skew part. For a
contracting flow, a rotation-invariant non-Gaussian $B$-self-decomposable law exists exactly when the
ratio stays below $d + 2$: an anisotropy budget of $2$, the ceiling of the stable index. The bound is
strict, unlike \ref{eq:stable-exponent-cone}, because the \Levy{} measure must integrate $|x|^2$ near
the origin. Clause (1) is proved in Appendix~\ref{app:wedge-axis}, clause (2) in
Appendix~\ref{app:proofs}.

\paragraph{Prior art.} The families of one linear flow are Urbanik's operator-self-decomposable laws
\sscite{urbanik1972levy}. Their exponents are Jurek's U-exponents with the sign reversed. Jurek proves
that the set of them is a closed convex cone and a Lie wedge, a cone of matrices closed under the
conjugations $e^{X}Ye^{-X}$ by its own edge, and computes the Gaussian's
(\sscite{jurek1992operator}, Prop.~1.2, Cor.~1.3 and the Example, pp.~1055--1057); for the standard
Gaussian it is $\{\mathrm{sym}\,B \ge 0\}$, the closure of the union of the cones that \ref{eq:stable-exponent-cone} defines for
$0 < \alpha < \infty$. The book of Jurek and Mason
restates this frame, and in its terms $\mathfrak k_\alpha$ is the set of exponents of $S_\alpha(I)$
(\sscite{jurek1993operator}, p.~85, Prop.~3.3.6, p.~86, and Prop.~3.3.10(a), p.~88). Their
Proposition~3.4.1 characterizes a contracting exponent $B$ of a law by its \Levy{} measure $\nu$, with
a matching condition on the Gaussian part (\sscite{jurek1993operator}, Prop.~3.4.1, p.~90). The
condition on $\nu$ is $\nu(U) \ge \nu(e^{tB}U)$ for every Borel set $U$ and every $t \ge 0$. For
contracting $B$ that reduces membership in $\mathfrak k_\alpha$ to a computation on the density $n_I$. The step set of an orbit corresponds to Urbanik's decomposability
semigroup of $\mu_I$, transported to the orbit (\sscite{jurek1993operator}, Lemma~2.3.3, p.~54). Jurek
and Mason compute that semigroup for Gaussian laws, where it gives the whole Loewner order
(Prop.~2.3.8, p.~58). The qualitative
fact behind Proposition~\ref{prop:anisotropy-budget} is Jurek's: a rotation-invariant law with such an
exponent outside the skew-symmetric matrices is self-decomposable (\sscite{jurek1992operator},
Prop.~1.6, p.~1056). Operator-stable laws, the limits of sums of independent copies normalized by matrices, have
exponents of their own, a different set
(\sscite{meerschaert2001limit}, Thm.~7.2.11, p.~262). Hudson and Mason compute them in the plane, and
Holmes, Hudson and Mason in every dimension, for the laws whose symmetry group is conjugate to the
orthogonal group (\sscite{hudson1981operator}, Thm.~3, p.~483; \sscite{holmes1982operator}, Thm.~3 and
Cor.~3, p.~603, and p.~606). For $S_\alpha(I)$ they are $\alpha^{-1}I$ plus the skew-symmetric
matrices. By Jurek's (1.7), read in the sign of this paper, the stability exponents of an infinitely
divisible law are exponents of decomposability (\sscite{jurek1992operator}, p.~1057), so these matrices
lie in
$\mathfrak k_\alpha$, and on them $\mathrm{sym}\,B$ is a multiple of $I$. None
of these sources computes the exponents of self-decomposability of a non-Gaussian law, bounds $\tr B$
by $\lambda_{\min}(\mathrm{sym}\,B)$, or compares two members of the stable orbit. So the cone
$\mathfrak k_\alpha$, the quantitative budget of Proposition~\ref{prop:anisotropy-budget} with its sharpness
through $S_\alpha(I)$, and the step cone of Proposition~\ref{prop:stable-step-cone} are new against the
sources read. The \Levy{} measure of the elliptical stable law, which the step cone compares, is
classical (\sscite{hudson1981operator}, Thm.~3, p.~483, in the plane; \sscite{holmes1982operator},
proof of Thm.~3, pp.~605--606, in every dimension).
% Q-0309 (fix group 5): Jurek 1992 and Hudson-Mason 1981 read at the page by Q-0308
% (records/affine/READING-jurek1992operator.md, READING-hudson1981operator.md; the table of
% records/affine/PRIOR-ART-novelty-claims.md). Jurek's U-exponent Q has mu = e^{tQ}mu * nu_t, so the
% exponent B here is -Q. Neither source covers K_alpha, the budget or the step cone; both are cited
% for the frame (Prop. 1.2, Cor. 1.3, the Example, (1.7)), the qualitative budget (Prop. 1.6) and the
% elliptical stable Levy measure (Hudson-Mason Thm. 3). The TODO(source) of section 1 is closed.
% Holmes-Hudson-Mason read at the page (Q-0294; records/affine/READING-holmes1982operator.md):
% stability exponents only (Thm. 1, p. 602; Cor. 3, p. 603); no exponents of self-decomposability,
% no anisotropy budget, no step cone. Their transform exp{i(a,y) - beta|Wy|^gamma} is written here in
% words (review flag 24: the source's letters a, beta, W, gamma carry other meanings in this paper).
% That alpha^{-1}I + Q satisfies eq:stable-exponent-cone is the line (d + alpha)/alpha >= d/alpha.

% CHANGED Q-0281 wave 2: blind review W2F-6 ("through" to "in the sense of": the proof of
% prop:stable-order-generated no longer goes through thm:cone-orders).
The step cone and the exponent cone are one object read in two ways. The chain of cones $C_\alpha$
of \S\ref{sec:invariant-wedges} generates, in the sense of Theorem~\ref{thm:cone-orders}, exactly the stable
orders, and its members of index $0 < \alpha < 2$ are the symmetric parts of the exponent cones. The link is the
\emph{tangent cone} of a closed semigroup $S$ of matrices, which the Lie theory of semigroups calls its
tangent wedge: the set of $B$ with $e^{tB} \in S$ for every $t \ge 0$, the generators of the
one-parameter semigroups that $S$ contains (\sscite{hilgert1989lie}, Prop.~V.1.7, pp.~375--376).

% shared with blueprint prop:stable-order-generated
\begin{proposition}[the stable step cone is the order of the cone field $C_\alpha$]\label{prop:stable-order-generated}
Let $d \ge 1$. For $0 \le \alpha < \infty$ put
\[
  C_\alpha := \{X \in \Sym(d) : X \text{ positive semidefinite},\ (d+\alpha)\,\lambda_{\min}(X) \ge \tr X\} ,
  \qquad C_\infty := \PSD(d) ,
\]
and write $\preceq_{C_\alpha}$ for the order of the cone field $C_\alpha$ (Theorem~\ref{thm:cone-orders}), $\lambda_{\min}$ the least
eigenvalue. Each $C_\alpha$ is an $\Orth(d)$-invariant closed convex cone containing $I$, with
$C_\alpha \cap (-C_\alpha) = \{0\}$. For $\Sigma, \Sigma' \in \PD(d)$ let
$M := \Sigma^{-1/2}\Sigma'\Sigma^{-1/2}$.
\begin{enumerate}
\item \emph{(Three orders, each generated by its cone field.)} $C_0 = \{cI : c \ge 0\}$, and
$\Sigma \preceq_{C_0} \Sigma'$ if and only if $\Sigma' = c\,\Sigma$ with $c \ge 1$, the ray order on $\PD(d)$.
For $0 < \alpha < \infty$ the condition $X \ge 0$ in $C_\alpha$ is implied by the inequality, and
$\Sigma \preceq_{C_\alpha} \Sigma'$ if and only if $\lambda_{\min}(M)^{d+\alpha} \ge \det M$; for
$0 < \alpha < 2$ this is the order $\preceq_\alpha$ of Proposition~\ref{prop:stable-step-cone}(2). $\Sigma \preceq_{C_\infty}
\Sigma'$ if and only if $\Sigma \le \Sigma'$, the Loewner order. In each case the related pairs are exactly
the pairs joined by a $C_\alpha$-causal curve, and $\Sigma^{1/2}M^t\Sigma^{1/2}$ is one.
\item \emph{(The exponent cone is the tangent cone.)} For $0 < \alpha < 2$ the set $\mathfrak k_\alpha$ of
\ref{eq:stable-exponent-cone} equals $\{B \in \gl(d) : e^{tB}e^{tB^{\tp}} \in \exp C_\alpha \text{ for every }
t \ge 0\}$, the tangent cone of the semigroup $\{A : I \preceq_{C_\alpha} AA^{\tp}\}$, the set of $B$ with $e^{tB}$ in
it for every $t \ge 0$, and
$\mathfrak k_\alpha \cap \Sym(d) = C_\alpha$.
\item \emph{(The Loewner order in other coordinates.)} For $0 < \alpha < \infty$ let
$\Phi_\alpha(\Sigma) := (\det\Sigma)^{-1/(d+\alpha)}\,\Sigma$ on $\PD(d)$. Then $\Phi_\alpha$ is a
diffeomorphism of $\PD(d)$ onto itself with inverse $T \mapsto (\det T)^{1/\alpha}\,T$;
$\Sigma \preceq_{C_\alpha} \Sigma'$ if and only if $\Phi_\alpha(\Sigma) \le \Phi_\alpha(\Sigma')$; the derivative of
$\Phi_\alpha$ at $\Sigma$ maps the cone $\Sigma^{1/2}C_\alpha\Sigma^{1/2}$ onto $\PSD(d)$; and
$\Phi_\alpha(A\Sigma A^{\tp}) = \psi_\alpha(A)\,\Phi_\alpha(\Sigma)\,\psi_\alpha(A)^{\tp}$ for
$A \in \GL(d)$, where $\psi_\alpha(A) := |\det A|^{-1/(d+\alpha)}A$ is an automorphism of the group
$\GL(d)$. In particular $\psi_\alpha$ maps the semigroup $\{A : I \preceq_{C_\alpha} AA^{\tp}\}$ onto the semigroup
$\{A : AA^{\tp} \ge I\}$ of the Loewner order.
\end{enumerate}
\end{proposition}

So up to a change of coordinates the stable orders are the Loewner order: $\preceq_\alpha$ is the
Loewner order read through $\Phi_\alpha$. The parameter interpolates between the rays, as $\alpha \to 0$,
where $\Phi_\alpha$ degenerates to the shape map $\Sigma \mapsto (\det\Sigma)^{-1/d}\Sigma$, and the
Loewner order at $\alpha = \infty$. The map $\Phi_\alpha$ acts on the parameters only.
The proof is in Appendix~\ref{app:wedge-axis}.


The cones form a chain, $C_0 \subsetneq C_\alpha \subsetneq C_\beta \subsetneq C_2 \subsetneq C_\infty$
for $0 < \alpha < \beta < 2$. For $\alpha < \infty$ every nonzero member of $C_\alpha$ is positive
definite, while $C_\infty$ contains the singular directions. Which general linear families each cone
carries is Theorem~\ref{thm:wedge-axis}, stated and proved in
\S\ref{sec:wedge-axis-proof}.

\subsubsection*{Every profile: the dilation ratio}

A rotation-invariant law with no Gaussian part and a \Levy{} density is described by one function of
the radius, its \emph{profile} $k$: the \Levy{} density at $x$ is $k(|x|)\,|x|^{-d}$, so that the
displacements with radius between $r_1$ and $r_2$ occur at the rate
$|S^{d-1}|\int_{r_1}^{r_2}k(r)\,dr/r$. It is the polar profile of
Definition~\ref{def:isotropic-origin-datum} divided by the area of the sphere, $k = k_d/|S^{d-1}|$. The
stable profile is $k(r) = c_{d,\alpha}^{-1}r^{-\alpha}$ (Lemma~\ref{lem:elliptical-stable}). For a
general nonincreasing profile the ordering of \Levy{} densities is read off one function, the
\emph{dilation ratio} $R(\lambda)$, the least factor by which the profile, the rate of displacements per
unit of $\log r$ in each direction, grows when the radius is divided by $\lambda$. It is $\lambda^\alpha$ for the stable profile and $1$ for a
profile bounded at the origin (Proposition~\ref{prop:orbit-step-cone}(3)).

% shared with blueprint prop:orbit-step-cone
\begin{proposition}[the ordered-density step set of an affine orbit]\label{prop:orbit-step-cone}
Let $d \ge 1$ and let $k : (0,\infty) \to [0,\infty)$ be nonincreasing, not identically $0$, and put
$m(r) := k(r)\,r^{-d}$, with $\int_{\RR^d}(|x|^2 \wedge 1)\,m(|x|)\,dx < \infty$. For $\lambda \ge 1$ put
\[
  R(\lambda) := \inf\{k(r/\lambda)/k(r) : r > 0,\ k(r) > 0\} \in [1, \infty) .
\]
\begin{enumerate}
\item \emph{(The orbit.)} There is a unique probability measure $\mu_I$ on $\RR^d$ with
$\widehat\mu_I(\omega) = e^{-\psi_I(\omega)}$, $\psi_I(\omega) := \int(1 - \cos\langle\omega,x\rangle)\,
m(|x|)\,dx$; it is rotation invariant and infinitely divisible with generating triplet
$(0, m(|x|)\,dx, 0)$. For $\Sigma \in \PD(d)$ its orbit $\mu_\Sigma := (\Sigma^{1/2})_*\mu_I$ has
$\widehat\mu_\Sigma(\omega) = \exp[-\int(1 - \cos\langle\omega,x\rangle)\,n_\Sigma(x)\,dx]$ with
$n_\Sigma(x) := (\det\Sigma)^{-1/2}m(|\Sigma^{-1/2}x|)$.
\item \emph{(The steps.)} For $\Sigma, \Sigma' \in \PD(d)$ put $M := \Sigma^{-1/2}\Sigma'\Sigma^{-1/2}$.
Then $n_\Sigma \le n_{\Sigma'}$ everywhere if and only if $\lambda_{\min}(M) \ge 1$ and
\[
  \det M \le \lambda_{\min}(M)^d\,R\bigl(\lambda_{\min}(M)^{1/2}\bigr)^2 ;
\]
and then $\Sigma \le \Sigma'$, and $\mu_{\Sigma'} = \mu_\Sigma * \rho$ with $\rho$ the infinitely
divisible probability measure with generating triplet $(0, (n_{\Sigma'} - n_\Sigma)(x)\,dx, 0)$.
\item \emph{(Which steps.)} The relation of (2) is a partial order on $\PD(d)$, invariant under
$\Sigma \mapsto A\Sigma A^{\tp}$ for $A \in \GL(d)$; it contains every ray $\Sigma \le \lambda^2\Sigma$,
$\lambda \ge 1$, and no step $\Sigma \to \Sigma + P$ with $0 < \rank P < d$. If
$k(r) = c\,r^{-\alpha}$, $0 < \alpha < 2$, then $R(\lambda) = \lambda^\alpha$ and its step set is the cone of
Proposition~\ref{prop:stable-step-cone}. If $k(0{+}) := \lim_{r\downarrow0}k(r)$ is finite, then $R \equiv 1$ and
its step set is the rays.
\end{enumerate}
\end{proposition}

The least eigenvalue of $M$ fixes how far a step may widen its narrowest direction, and $R$ how much
volume that buys. By the polar criterion of self-decomposability, C5 (\S\ref{sec:cited-facts};
\S\ref{sec:iso-recap}), the laws of clause (1) are the rotation-invariant
self-decomposable laws with no Gaussian part whose \Levy{} measure has a density, so clause (2)
applies to the orbit of every such law. The proposition gives sufficiency. That a positive increment forces ordered densities is proved for the stable
profiles (Lemma~\ref{lem:stable-tail}) and for the profiles bounded at the origin
(Proposition~\ref{prop:bounded-profile-rays}). For a general profile it is open: the positivity of an
increment alone does not exclude a signed \Levy{} measure. Inside a cascade family along a stable order or
a cone order it holds for every profile, since there the increments are infinitely divisible
(Proposition~\ref{prop:wedge-axis-profiles}). The class $I_0$ does not close the gap. By
Cuppens's criterion, a law of clause (1) whose \Levy{} density is continuous and positive off the
origin lies outside the class $I_0$ of laws whose factors are all infinitely divisible
(\sscite{cuppens1975decomposition}, Thm.~7.2.1, p.~157). The stable and \Matern{} laws are among them
(for the stable laws on the line, \sscite{linnik1977decomposition}, Thm.~6.7.2 and Ex.~1, p.~218).

The proof substitutes $x = \Sigma^{1/2}y$, which turns the ordering of densities into a comparison of
$m$ at $|y|$ and at $\langle y,M^{-1}y\rangle^{1/2}$. On each sphere the comparison is hardest along an
eigenvector for $\lambda_{\min}(M)$, where it reads $k(r/\lambda) \ge (\det M)^{1/2}\lambda^{-d}k(r)$ with
$\lambda := \lambda_{\min}(M)^{1/2}$, and the definition of $R$ makes this the inequality of (2). The
complete proof is in Appendix~\ref{app:wedge-axis}.

% Round 3 (report D, item 5; VERIFY-round3-landscape-orbits-paper.md section 9; taken unless vetoed,
% PLAN-review-response-paper.md R3.1): the dyadic counterexample, as in the annotation of
% prop:orbit-step-cone; an unnumbered paragraph, so that no later statement is renumbered.
\paragraph{A step set, not a cone.} The relative steps $M$ allowed by clause (2) form a set
$\mathcal K_R$ invariant under orthogonal conjugation, a step set in the sense of
Proposition~\ref{prop:invariant-orders}. For the stable profiles and the profiles bounded at the origin
it is $\exp C$ for a cone $C$ (clause (3)), and in general it need not be. The dyadic profile of
Lemma~\ref{lem:dilation-ratio}(5) below has $R = 1$ on $[1, 2)$ and $R(2) = 2$. For $d \ge 2$, the
matrix $M := \diag(4, \dots, 4, 16)$ has $\det M = 4^d\cdot 4 = \lambda_{\min}(M)^dR(2)^2$, so
$M \in \mathcal K_R$. Its square root $M^{1/2} = \diag(2, \dots, 2, 4)$ has $\det M^{1/2} = 2^d\cdot 2$,
above $\lambda_{\min}(M^{1/2})^dR(\sqrt2)^2 = 2^d$, so $M^{1/2} \notin \mathcal K_R$. A set $\exp C$ contains
$e^{X/2}$ with $e^X$, so this $\mathcal K_R$ is $\exp C$ for no cone $C$.

The dilation ratio is monotone and supermultiplicative, and it has a growth exponent. The growth
exponent bounds $R$ from above, so the rung that ordered densities give a profile is its lower
exponent instead.

% shared with blueprint lem:dilation-ratio+4c533f713cd8
\begin{lemma}[the dilation ratio: monotonicity and two exponents]\label{lem:dilation-ratio}
Let $d \ge 1$, let $k$ be a profile as in Proposition~\ref{prop:orbit-step-cone}, and let $R$ be its dilation ratio,
$R(\lambda) = \inf\{k(r/\lambda)/k(r) : r > 0,\ k(r) > 0\}$ for $\lambda \ge 1$.
\begin{enumerate}
\item \emph{(Monotone and supermultiplicative.)} $R(1) = 1$, $R$ is nondecreasing on $[1,\infty)$, and
$R(\lambda\nu) \ge R(\lambda)\,R(\nu)$ for all $\lambda, \nu \ge 1$.
\item \emph{(Below the square.)} $R(\lambda) < \lambda^2$ for every $\lambda > 1$.
\item \emph{(Two exponents.)} Put $\beta_-(k) := \inf_{\lambda>1}\log R(\lambda)/\log\lambda$. The
limit $\beta(k) := \lim_{\lambda\to\infty}\log R(\lambda)/\log\lambda$ exists and equals
$\sup_{\lambda>1}\log R(\lambda)/\log\lambda$. They satisfy $0 \le \beta_-(k) \le \beta(k) \le 2$ and
$\beta_-(k) < 2$, and
\[
  \lambda^{\beta_-(k)} \le R(\lambda) \le \lambda^{\beta(k)} \qquad (\lambda \ge 1) .
\]
For $\alpha \ge 0$, $R(\lambda) \ge \lambda^\alpha$ for every $\lambda \ge 1$ if and only if
$\alpha \le \beta_-(k)$.
\item \emph{(The rung.)} Let $d \ge 2$ and $0 < \alpha < 2$, and let $\preceq_\alpha$ be the order of
Proposition~\ref{prop:stable-step-cone}(2). Every pair $\Sigma \preceq_\alpha \Sigma'$ satisfies $n_\Sigma \le n_{\Sigma'}$
everywhere, in the notation of Proposition~\ref{prop:orbit-step-cone}(1), if and only if $\alpha \le \beta_-(k)$.
\item \emph{(Values.)} If $k(r) = c\,r^{-\alpha}$ with $c > 0$ and $0 < \alpha < 2$, then
$\beta_-(k) = \beta(k) = \alpha$. If $k(0{+})$ is finite, then $\beta_-(k) = \beta(k) = 0$. The profile
$k(r) := 2^n$ for $2^{-n-1} < r \le 2^{-n}$, $n \in \ZZ$, satisfies the hypotheses of
Proposition~\ref{prop:orbit-step-cone} and has $\beta_-(k) = 0$ and $\beta(k) = 1$.
\end{enumerate}
\end{lemma}

So the condition $R(\lambda) \ge \lambda^\alpha$ for every $\lambda \ge 1$, which appears in
Theorem~\ref{thm:wedge-axis}(3) and Proposition~\ref{prop:similarity-stable-order-families}(2), is
$\alpha \le \beta_-(k)$, and for $d \ge 2$ it is exactly the condition that the ordered-density step set of
$k$ contain the stable order of index $\alpha$. The two exponents agree for the stable profiles, for
the profiles bounded at the origin and for the tempered profiles below, and the
dyadic profile of clause (5) shows that they need not agree in general. The value $\beta(k) = 2$
occurs, for $k(r) = r^{-2}(1 + \log(1/r))^{-2}$ on $(0, 1]$ and $r^{-2}$ beyond. The lower exponent
also decides which stable orders the orbit's family survives, for every profile and with
any Gaussian part (Proposition~\ref{prop:wedge-axis-profiles}). Whether a single positive increment outside a family orders
the densities is the open converse of Proposition~\ref{prop:orbit-step-cone}. The proof is in
Appendix~\ref{app:proofs}.

\paragraph{Profiles stable at the origin.} A profile $k(r) = r^{-\alpha}j(r)$, with $j$ nonincreasing
and $j(0{+})$ finite and positive, has the dilation ratio $\lambda^\alpha$ of the stable profile. So
its orbit's ordered-density step set is the stable cone of index $\alpha$, however the profile is tempered
at large radii. With an exponential tempering the orbit has exponential moments, as the \Matern{} law
does and the stable law does not. One such family has the transforms
$\exp[-((1 + \langle\omega,\Sigma\omega\rangle)^{\alpha/2} - 1)]$. With its exponent normalized to
$(2\gamma/\alpha)((1 + u)^{\alpha/2} - 1)$, its limit as $\alpha \downarrow 0$ is the \Matern{} orbit of
index $\gamma$, whose profile is bounded at the origin. These are Proposition~\ref{prop:tempered-step-cone},
in Appendix~\ref{app:tempered}.

\subsubsection*{Index $0$: the rays}

At the bottom of the ladder, rung $0$, are the orbits whose transform falls off like a power of the
frequency: $e^{-g(cu)}/e^{-g(u)} \to c^{-\kappa}$ as $u \to \infty$, for every $c > 0$. The \Matern{}
orbit of index $\gamma$ is one, with $\kappa = \gamma$. For them one positive increment between two
apertures forces a ray. The statement is about a single pair of apertures of the orbit, so it holds
along every order.

% shared with blueprint prop:log-orbit-rays+739cc86e19da
\begin{proposition}[orbits of logarithmic exponent cascade along rays only]\label{prop:log-orbit-rays}
Let $d \ge 1$, let $g$ satisfy the hypothesis of Proposition~\ref{prop:affine-orbits}, and let $(\mu_\Sigma)$
be its affine orbit, $\widehat\mu_\Sigma(\omega) = e^{-g(\langle\omega,\Sigma\omega\rangle)}$. Suppose
there is $\kappa > 0$ with
\[
  g(cu) - g(u) \to \kappa\log c \qquad (u \to \infty), \text{ for every } c > 0 .
\]
If $\Sigma, \Sigma' \in \PD(d)$ and $\mu_{\Sigma'} = \mu_\Sigma * \rho$ for a positive finite measure
$\rho$, then $\Sigma' = c\,\Sigma$ for some $c > 0$. In particular, for every $\gamma > 0$, no step
of the \Matern{} orbit of index $\gamma$, $\widehat\mu_\Sigma(\omega) = (1 + \langle\omega,\Sigma
\omega\rangle)^{-\gamma}$ (Proposition~\ref{prop:affine-orbits}(4)), off the rays has a positive increment.
\end{proposition}

At small scales the increment looks like the distribution whose transform is
$(\langle\omega,\Sigma\omega\rangle/\langle\omega,\Sigma'\omega\rangle)^\kappa$, a function homogeneous of
degree $0$. A positive measure that is invariant under dilations and finite on compact sets is a point
mass, so that transform must be constant.
Which rays carry a positive increment is a separate question: for a general orbit, exactly those of
a self-decomposable $\mu_I$ (Proposition~\ref{prop:ray-families}).
The proof is in Appendix~\ref{app:proofs}.

\begin{remark}[the analytic step in the literature]\label{rem:log-orbit-crum}
The proof in Appendix~\ref{app:proofs} shows that the limit
$h(\omega) := \lim_{t\to\infty}\widehat\rho(t\omega)$, which is
$(\langle\omega,\Sigma\omega\rangle/\langle\omega,\Sigma'\omega\rangle)^\kappa$ for $\omega \ne 0$, is
constant off $0$. That step has a second, published route. The limit $h$ of $\widehat\rho(t\,\cdot)$, with
$h(0) = 1$, is a pointwise limit of continuous functions, so it is measurable. As a pointwise limit of
characteristic functions it is positive definite (\sscite{sasvari2013multivariate}, Thm.~1.4.10,
p.~28). A measurable positive definite
function on $\RR^d$ equals a continuous one $f_c$ almost everywhere (\sscite{crum1956positive},
Theorem, p.~548; \sscite{sasvari2013multivariate}, Thm.~2.11.7, p.~130). As $h$ is continuous off
$0$, $h = f_c$ off $0$. So $f_c$ is homogeneous of degree $0$ off $0$ and continuous at $0$, hence
constant, and $h$ is that constant off $0$.
\end{remark}

For the \Matern{} orbit of every index the two halves combine: the rays carry positive increments,
and nothing else does.

% shared with blueprint prop:matern-steps+2810f2866dc2
\begin{proposition}[the \Matern{} orbit of every index: the positive steps are the rays]\label{prop:matern-steps}
Let $d \ge 1$, $\gamma > 0$, and let $(\mu_\Sigma)_{\Sigma\in\PD(d)}$ be the \Matern{} orbit of index
$\gamma$, $\widehat\mu_\Sigma(\omega) = (1 + \langle\omega,\Sigma\omega\rangle)^{-\gamma}$
(Proposition~\ref{prop:affine-orbits}(4)). Let $\Sigma, \Sigma' \in \PD(d)$ with $\Sigma' \ne \Sigma$. Then there
is a positive finite measure $\rho$ with $\mu_{\Sigma'} = \mu_\Sigma * \rho$ if and only if
$\Sigma' = \lambda^2\Sigma$ for some $\lambda > 1$; and then $\rho$ is unique, and it is the
probability measure $(\Sigma^{1/2})_*\nu_\lambda$, with $\nu_\lambda$ the law of $\sqrt{2Y}\,Z$ of
Proposition~\ref{prop:affine-orbits}(4).
\end{proposition}

No hypothesis $\Sigma \le \Sigma'$ is needed: no pair of distinct apertures off a ray is joined by a
positive increment, comparable or not.

% proof of prop:matern-steps
\begin{proof}
$\widehat\mu_\Sigma$ has no zero, so an increment $\rho$, if there is one, has the transform
$\widehat\rho(\omega) = \bigl((1 + \langle\omega,\Sigma\omega\rangle)/(1 + \langle\omega,\Sigma'\omega\rangle)\bigr)^\gamma$,
which determines it. \emph{Only if.} By Proposition~\ref{prop:log-orbit-rays}, $\Sigma' = c\,\Sigma$ for
some $c > 0$, and $c \ne 1$ as $\Sigma' \ne \Sigma$. As $\widehat\rho(0) = 1$, $\rho$ is a probability
measure and $|\widehat\rho| \le 1$; at $\omega \ne 0$, with $u := \langle\omega,\Sigma\omega\rangle > 0$,
$\widehat\rho(\omega) = ((1+u)/(1+cu))^\gamma$, which exceeds $1$ if $c < 1$. So $c > 1$, and
$\lambda := c^{1/2}$. \emph{If.} By Proposition~\ref{prop:affine-orbits}(4), the hypothesis of
Proposition~\ref{prop:affine-orbits}(2) holds for $\lambda$ with the stated $\nu_\lambda$, and
Proposition~\ref{prop:affine-orbits}(2) gives $\mu_{\lambda^2\Sigma} = \mu_\Sigma * (\Sigma^{1/2})_*\nu_\lambda$.
\end{proof}

A filter computed by finitely many elliptic solves in the shape's own operator has a rational symbol,
and a decaying rational symbol has a logarithmic exponent.

% shared with blueprint prop:rational-symbols
\begin{proposition}[rational symbols cascade along rays only]\label{prop:rational-symbols}
Let $d \ge 1$, and let $P$ and $Q$ be real polynomials with $\operatorname{deg}P > \operatorname{deg}Q$ and $P(0) = Q(0)$,
neither of which vanishes on $[0,\infty)$. Suppose that $\omega \mapsto Q(|\omega|^2)/P(|\omega|^2)$
is the characteristic function of a probability measure $\mu_I$ on $\RR^d$, and let
$\mu_\Sigma := (\Sigma^{1/2})_*\mu_I$ for $\Sigma \in \PD(d)$, so that
$\widehat\mu_\Sigma(\omega) = Q(\langle\omega,\Sigma\omega\rangle)/P(\langle\omega,\Sigma\omega\rangle)$.
Then $g(u) := \log(P(u)/Q(u))$ satisfies the hypothesis of Proposition~\ref{prop:affine-orbits}, and that of
Proposition~\ref{prop:log-orbit-rays} with $\kappa = \operatorname{deg}P - \operatorname{deg}Q$. So if $\Sigma, \Sigma' \in \PD(d)$ and
$\mu_{\Sigma'} = \mu_\Sigma * \rho$ for a positive finite measure $\rho$, then $\Sigma' = c\,\Sigma$ for
some $c \ge 1$. This applies, with $n \ge 1$ and $a_1, \dots, a_n > 0$, to the \Matern{} orbits of integer index
$\gamma$ ($Q = 1$, $P = (1+u)^\gamma$, $\kappa = \gamma$), to the products
$\prod_{i\le n}(1 + a_iu)^{-1}$ ($\kappa = n$), and to the mixtures $\sum_{i\le n}w_i(1 + a_iu)^{-1}$
with $w_i \ge 0$ and $\sum_iw_i = 1$ ($\kappa = 1$).
\end{proposition}

Write $\mathcal L_\Sigma := -\nabla\cdot\Sigma\nabla$, the operator with symbol $u = \langle\omega,\Sigma\omega\rangle$.
Finitely many solves of shifted equations $(\mathcal L_\Sigma - z)v = w$, $z \in \CC \setminus [0,\infty)$, and
finitely many applications of $\mathcal L_\Sigma$ multiply the transform by a rational function of $u$.
Conversely $Q(u)/P(u)$, with $P$ factored over $\CC$, is computed by $\operatorname{deg}P$ such solves
and one application of $Q(\mathcal L_\Sigma)$. A decaying rational symbol behaves like $u^{-\kappa}$ at infinity
with $\kappa$ a positive integer. So every affine orbit of a law that such solves implement exactly,
with a transform that does not vanish and decays, changes shape only by re-measurement. The converse
fails: the \Matern{} orbits of every index and the orbits of profiles bounded at the origin cascade
along rays only too (Propositions~\ref{prop:matern-steps} and~\ref{prop:bounded-profile-rays}), and the
tempered orbits of Proposition~\ref{prop:tempered-step-cone}(3), whose symbols are not rational, have
at least the stable cone. Symbols that do not decay, with an atom at the origin, are not covered.
The proof is in Appendix~\ref{app:proofs}.


The profiles bounded at the origin close the converse of Proposition~\ref{prop:orbit-step-cone}(2):
by Lemma~\ref{lem:isotropic-exponent-log}(2) such a profile gives a logarithmic exponent.

% shared with blueprint prop:bounded-profile-rays+ec7f6de3b73c
\begin{proposition}[orbits of a profile bounded at the origin cascade along rays only]\label{prop:bounded-profile-rays}
Let $d \ge 1$, let $k$ satisfy the hypotheses of Proposition~\ref{prop:orbit-step-cone} with
$k(0{+}) := \lim_{r\downarrow0}k(r)$ finite, and let $(\mu_\Sigma)$ be the orbit of
Proposition~\ref{prop:orbit-step-cone}(1), with \Levy{} densities $n_\Sigma$. For a unit vector $e$ the function
$g(u) := \psi_I(\sqrt u\,e)$ satisfies the hypothesis of Proposition~\ref{prop:affine-orbits}, and that of
Proposition~\ref{prop:log-orbit-rays} with $\kappa = \tfrac12|S^{d-1}|\,k(0{+}) > 0$. Consequently, for
$\Sigma, \Sigma' \in \PD(d)$ the following are equivalent:
\begin{enumerate}
\item $\mu_{\Sigma'} = \mu_\Sigma * \rho$ for a positive finite measure $\rho$;
\item $n_\Sigma \le n_{\Sigma'}$ everywhere;
\item $\Sigma' = c\,\Sigma$ for some $c \ge 1$.
\end{enumerate}
\end{proposition}

So for every profile bounded at the origin a positive increment, ordered \Levy{} densities and a ray
are the same condition. The class contains the \Matern{} orbits
(Proposition~\ref{prop:tempered-step-cone}(4), with $\kappa = \gamma$), every subordinated isotropic
datum with $a = 0$ (Definition~\ref{def:isotropic-origin-datum}) whose delay profile is nonzero and bounded at the
origin, and the unit-step profiles
$k = \mathbf 1_{(0,\tau)}$. The equivalence holds on both sides of the boundedness threshold of the isotropic paper
\prog{Iso}{Thm.~3.13}, the value $k_d(0{+}) = d$ of the polar profile at the origin above which a
kernel without Gaussian part is bounded at its centre and at or below which it is not: every bounded
profile has $R \equiv 1$, whatever its value at the origin. For a profile unbounded at the origin the
three conditions part: for the stable profiles (1) and (2) are equivalent and (3) fails
(Proposition~\ref{prop:stable-step-cone}), and for the others it is open whether (1) implies (2).
The proof is in Appendix~\ref{app:proofs}.

% Round 2 (presentation review), item 11 (VERIFY-round2-notation-paper.md § 11): the mechanism with
% its formula and the hypotheses of Proposition E.4 (d >= 2, Sigma in PD(d), 0 != P in PSD(d));
% compared with Remark E.3 and Proposition E.4(1)-(3), with B_d = 1/|S^{d-1}|.
\paragraph{The size of the failure.} For the \Matern{} orbit of index $1$, $d \ge 2$, and a step from
$\Sigma \in \PD(d)$ to $\Sigma' = \Sigma + P$ with $0 \ne P \in \PSD(d)$, the would-be increment has the
transform $1 - \langle\omega,P\omega\rangle\,\widehat\mu_{\Sigma'}(\omega)$. Off the origin it is the
function $q = \nabla\cdot P\nabla K_{\Sigma'}$, where the density $K_{\Sigma'}$ of $\mu_{\Sigma'}$ is a
linear image of the Green's function of $1 - \Delta$. Near the origin, in the coordinates
$y = \Sigma'^{-1/2}x$ and with $P' := \Sigma'^{-1/2}P\Sigma'^{-1/2}$, $q$ is
$(\det\Sigma')^{-1/2}|S^{d-1}|^{-1}|y|^{-d}(d\langle\widehat y,P'\widehat y\rangle - \tr P')$ to leading
order. Unless $\Sigma'$ is a multiple of $\Sigma$, the bracket is negative around a least eigenvector
of $P'$. The negative mass on $\{a < |y| < 1\}$ then grows like $\log(1/a)$, so the increment is not
even a finite signed measure
(Proposition~\ref{prop:matern-cross-shape} with Remark~\ref{rem:matern-mechanism}, and
Proposition~\ref{prop:matern-negative-mass}, with a computation on a grid in
Remark~\ref{rem:matern-grid}; all four in Appendix~\ref{app:matern}).

\subsubsection*{The stable order on the positive semidefinite matrices}

The theorem on the general linear row (\S\ref{sec:wedge-axis-proof}) reads the families along the
stable order, extended to $\PSD(d)$ by the origin and the rays. The next proposition gives that order
and the general linear cell along it, on which the theorem rests. The same order for the similarities
is the subject of Proposition~\ref{prop:similarity-stable-order-kernels} in \S\ref{sec:smaller-groups}.

% Round-2 review (Gemini), item 9, the author's decision D8 (2026-11-03): the node is split; it keeps
% clauses (1)-(2), and clauses (3)-(5) are prop:similarity-stable-order-kernels, printed in section 4.9
% (05-smaller-groups.tex) with the lead sentences on eccentricity and spread and the comments on those
% clauses that followed here.
% shared with blueprint prop:similarity-stable-order-families+08facbc24d81
\begin{proposition}[the stable order on $\PSD(d)$ and the general linear row along it]\label{prop:similarity-stable-order-families}
Let $d \ge 2$ and $0 < \alpha < 2$. Write $\preceq_\alpha$ for the order of
Proposition~\ref{prop:stable-step-cone}(2) on $\PD(d)$, and $\operatorname{ecc}(\Sigma) :=
\lambda_{\max}(\Sigma)/\lambda_{\min}(\Sigma)$ for $\Sigma \in \PD(d)$.
\begin{enumerate}
\item \emph{(The stable order on $\PSD(d)$.)} For $\Sigma, \Sigma' \in \PSD(d)$ put
$\Sigma \le_\alpha \Sigma'$ when $\Sigma = 0$, or $\Sigma' = c\,\Sigma$ for some $c \ge 1$, or
$\Sigma, \Sigma' \in \PD(d)$ and $\Sigma \preceq_\alpha \Sigma'$. This is a partial order on $\PSD(d)$,
invariant under $\GL(d)$, containing the rays and contained in the Loewner order, and every pair
$\Sigma \le_\alpha \Sigma'$ is joined by a chain as in Theorem~\ref{thm:cascade-infinitely-divisible}(2); so
$(G, \le_\alpha)$ is a wedge datum with the path condition (P) for every closed subgroup $G$ of $\GL(d)$
containing the dilations. If $\alpha \le \beta < 2$, then $\Sigma \le_\alpha \Sigma'$ implies
$\Sigma \le_\beta \Sigma'$. For $\Sigma \in \PD(d)$ with least and largest eigenvalues
$\lambda_1 \le \lambda_d$, and $\varepsilon, t \ge 0$,
\[
  \varepsilon I \le_\alpha \Sigma \Leftrightarrow \varepsilon \le \varepsilon_\alpha(\Sigma) :=
  \Bigl(\frac{\lambda_1^{d+\alpha}}{\det\Sigma}\Bigr)^{1/\alpha} ,
  \qquad
  \Sigma \le_\alpha tI \Leftrightarrow t \ge t_\alpha(\Sigma) :=
  \Bigl(\frac{\lambda_d^{d+\alpha}}{\det\Sigma}\Bigr)^{1/\alpha} ,
\]
and $t_\alpha(\Sigma)/\varepsilon_\alpha(\Sigma) = \operatorname{ecc}(\Sigma)^{(d+\alpha)/\alpha}$.
\item \emph{(The general linear row.)} Let $k$ be a profile as in Proposition~\ref{prop:orbit-step-cone} whose
dilation ratio has $R(\lambda) \ge \lambda^\alpha$ for every $\lambda \ge 1$, and $\mu_I$ the law of that
proposition's clause (1). Then $\Sigma \mapsto (\Sigma^{1/2})_*\mu_I$, $\Sigma \in \PSD(d)$, is the family of
kernels from scale zero of a $(\GL(d), \le_\alpha)$-cascade family of kernels. This applies to the stable
law $S_\beta(I)$ of every index $\beta \in [\alpha, 2)$ (the profile $k(r) = c_{d,\beta}^{-1}r^{-\beta}$ of
Lemma~\ref{lem:elliptical-stable}), and to the profiles of Proposition~\ref{prop:tempered-step-cone} of every index
$\beta \in [\alpha, 2)$.
\end{enumerate}
\end{proposition}

The order $\le_\alpha$ of clause (1) is $\preceq^+_{C_\alpha}$ of Theorem~\ref{thm:cone-orders}(4), and the
general linear row along it is Theorem~\ref{thm:wedge-axis}(3). The thresholds of clause (1) place every
definite $\Sigma$ between two multiples of $I$ in the order, the sandwich on which the bounds of
Proposition~\ref{prop:similarity-stable-order-kernels} rest. The proof is in Appendix~\ref{app:wedge-axis}.
